\documentclass{article}
\usepackage[a4paper,margin=2.5cm]{geometry}
\usepackage{alltt}
\usepackage[hidelinks]{hyperref}
\usepackage{stickyflow}

\title{The stickyflow package}
\author{Stickyflow}
\date{Version 0.2.0, 26 November 2026}

% Listings show the environment with \textbackslash{}begin, so the package does not take them for a figure.
\newenvironment{code}{\par\smallskip\begin{alltt}\small}{\end{alltt}\smallskip}

\begin{document}
\maketitle

\section{What it does}

The stickyflow package draws figures written in SFL while a document compiles with LuaLaTeX. SFL is a language that describes a visual
artifact in text: its elements, what they mean, how they relate, where they sit and how they look. A chart, a diagram, a plan or an
interface is one block of SFL, and the package turns that block into a figure of the document.

The figure is drawn by the renderer of stickyflow.js itself, translated to Lua, so it looks the same in the PDF as in a browser. Its texts
are set by TeX in the fonts of the document. The PDF also carries the SFL of every figure, its data and its place on the page, so a viewer
that knows SFL can make the figure live again.

\begin{center}
\begin{sfl}[0.6\linewidth]
sfl 1.6

Model board {
  Chart skills {
    type: radar
    text: "Team profile"
    areaOpacity: 0.22
    categories: ["Speed", "Quality", "Cost", "Reliability", "Support", "Reach"]

    data {
      Axis value { min: 0; max: 10; step: 2 }

      now { label: "Today"; values: [8, 7, 5, 9, 6, 7]; color: "#0d9488" }
      goal { label: "Target"; values: [9, 9, 7, 8, 9, 8]; color: "#f4694c" }
    }

    style {
      radarStrokeWidth: 2.5
      radarPointSize: 4
      gridColor: "#e5e7eb"
      textColor: "#111827"
      showLegend: true
      cornerRadius: 14
      strokeWidth: 1
      stroke: "#e5e7eb"
      fill: "#ffffff"
    }
    layout {
      position: 0, 0
      size: 560, 440
    }
  }
}
\end{sfl}
\end{center}

The figure above is written in the source of this documentation.

\section{Requirements}

The package needs LuaLaTeX. It does not need shell escape and runs no program outside TeX, so it also works on Overleaf, where the compiler
is chosen under Menu, Compiler, LuaLaTeX. It loads \texttt{iftex}, \texttt{xcolor} and \texttt{graphicx}.

The package consists of two files, \texttt{stickyflow.sty} and \texttt{stickyflow.lua}. Where the package is not installed with the TeX
distribution, put both next to the document, or upload them into the Overleaf project.

\section{Figures in the document}

Load the package and write the SFL of a figure between \texttt{\textbackslash{}begin\{sfl\}} and \texttt{\textbackslash{}end\{sfl\}}:

\begin{code}
\textbackslash{}usepackage\{stickyflow\}
...
\textbackslash{}begin\{sfl\}[0.8\textbackslash{}linewidth]
sfl 1.6

Model board \{
  ...
\}
\textbackslash{}end\{sfl\}
\end{code}

The optional argument is the width of the figure. Without it the figure is as wide as the line. The figure keeps its proportions, so its
height follows from the width.

The lines of the environment are read before TeX sees them. Braces, hash signs and percent signs therefore stay what they are in SFL and
need no escaping. For the same reason the environment cannot stand inside the argument of another command. Inside a \texttt{figure} or a
\texttt{center} environment it is fine.

\section{Data}

A figure with data sources gets its rows as JSON after a line of three dashes:

\begin{code}
\textbackslash{}begin\{sfl\}[0.8\textbackslash{}linewidth]
sfl 1.10

DataSource sales \{
  mode: table
  fields \{ quarter: string; amount: number \}
\}
...
---
\{ "sales": [\{ "quarter": "Q1", "amount": 42 \}, \{ "quarter": "Q2", "amount": 55 \}] \}
\textbackslash{}end\{sfl\}
\end{code}

The figure is drawn with these rows and with the values its parameters start with.

\section{Figures from files}

A figure can stay in a file of its own, with its data in another:

\begin{code}
\textbackslash{}sflinput[0.62\textbackslash{}linewidth]\{section-study.sfl\}
\textbackslash{}sflinput[0.8\textbackslash{}linewidth][revenue.json]\{revenue.sfl\}
\end{code}

The first optional argument is the width, the second the file with the rows of the data sources. Unlike the environment,
\texttt{\textbackslash{}sflinput} can stand anywhere a box can.

\section{Scenes in three dimensions}

Only a browser draws a scene in three dimensions, so in print such a scene is a still picture set into its place. Give the picture as the
third optional argument, or put a picture with the name of the SFL file next to it, such as \texttt{robot-line.png} for
\texttt{robot-line.sfl}:

\begin{code}
\textbackslash{}sflinput[0.9\textbackslash{}linewidth]\{robot-line.sfl\}
\textbackslash{}sflinput[0.9\textbackslash{}linewidth][][still.png]\{robot-line.sfl\}
\textbackslash{}begin\{sfl\}[0.9\textbackslash{}linewidth][][still.png] ... \textbackslash{}end\{sfl\}
\end{code}

The models a scene names are attached to the PDF. The package finds a model by its path, or by the last part of the path, so
\texttt{/models/robot.glb} is also found as \texttt{robot.glb} next to the document.

\section{Figures with shaders}

A shader that fills a shape, such as a moving sea or a glowing gauge, is drawn by the graphics hardware of a browser as well, and TeX has
none. Print shows the color beneath the shader. When the look of a figure depends on its shaders, give it a picture as the third optional
argument, as for a scene. A figure without a scene in three dimensions is then that picture in print, at the size of the figure, while its
source still travels with the page and a viewer draws it live:

\begin{code}
\textbackslash{}sflinput[0.3\textbackslash{}linewidth][battery.json][battery.png]\{battery.sfl\}
\end{code}

\section{A sample paper}

The folder \texttt{examples/paper} holds a sample paper in two columns, with listings, equations and a table, whose seven figures show
what the package draws: a model fitted by hand to measurements, a field of densities and a hall of robots in three dimensions, a house and
a district drawn from their parameters, two instruments with shaders and a synthesizer. It compiles as it is, and its PDF lies beside it.

\section{Live figures}

After the last page the package attaches a file \texttt{stickyflow.json} to the PDF. It holds the SFL, the data and the place on the page of
every figure. A viewer that reads it can draw the figures again and let them react: sliders move, and what is bound to them follows.
Controls such as sliders are left out of print.

\section{Limits}

Gradients are drawn in their middle color. Shaders and input fields, which only a browser draws, are left out. Only the first scene in
three dimensions of a figure gets the still picture; further ones are a frame in print.

\section{How it is made}

The static core of stickyflow.js, which renders SFL to SVG without a browser, is translated to Lua with TypeScriptToLua.
\texttt{stickyflow.lua} holds that core and the part that draws the SVG as PDF operators: paths, shapes, arcs, arrow heads, clips, dashes
and transparency. Shapes go into forms with their own resources, so the resources of the page stay as other packages set them. The sources
and the build are at \url{https://github.com/Stickyflow/stickyflow-latex}.

\section{License}

The package is licensed under the Apache License 2.0. The translated core comes from stickyflow.js, which is licensed under the same
license.

\end{document}
